Le cahier des charges est un document qui définit les objectifs, les besoins, les livrables, les contraintes ainsi que les critères de validation d'un projet.  

Dans cette partie, nous évoquerons les éléments les plus importants, c'est-à-dire les objectifs, les données à disposition ainsi que les fonctionnalités attendues pour les deux visualisations retenues.

\subsection{Les objectifs}

Dans le cadre de notre projet, nous allons réaliser deux preuves de concept de visualisations pour l'application \gls{aikon}.

Pour les deux visualisations, nous nous allons nous appuyer sur les \textit{regions} de deux \textit{witnesses}.
La première servira à comparer les \textit{region pairs} de \textit{witnesses} issus du projet \gls{vhs}. La deuxième aura pour but de visualiser les \textit{regions} similaires de données extraites du projet \textit{eida} afin de comparer les résultats de la reconnaissance automatique de similarités avant l'annotation manuelle des chercheurs. 

\subsection{Les contraintes techniques}

Pour mener à bien ce projet, nous nous heurtons à quelques contraintes techniques. 

Tout d'abord, le temps de réalisation est limité à la durée du stage qui est de quatre mois.

Ensuite, il n'est pas possible de se connecter à la base de données car il s'agit de preuves de concept qui ne seront pas intégrées telles quelles dans l'application. 

Il faut aussi prendre en compte la qualité des données. Pour celles du projet \gls{vhs}, les résultats du traitement automatique ont été annotés par les chercheurs. Même s'il peut rester une part de subjectivité pour définir si deux \textit{regions} sont réellement des \textit{region pairs}, elles sont fiables. Les données issues du projet \gls{eida} sont les résultats bruts de la reconnaissance automatique des similarités. Il faut donc prêter attention à la fiabilité des \textit{region pairs}. De plus, il risque d'y avoir une quantité plus importante de données que si elles avaient été annotées au préalable par les chercheurs comme pour celles de \gls{vhs}. 

Pour finir, les visualisations doivent être adaptables à tous les projets utilisant l'application \gls{aikon}.

\subsection{Les données}

Etant donné que nous ne pouvons pas nous connecter directement à la base de données, nous allons donc utiliser des exports \gls{csv} de la table \textit{region-pairs}.

Pour la visualisation en \textit{bipartite graph}, nous utiliserons les données du projet \gls{vhs} et pour celle en \textit{force-directed layout}, celles du projet \gls{eida}.

\subsection{Les fonctionnalités}

\subsubsection{Bipartite graph}

Dans la visualisation \textit{Bipartite graph}, les \textit{regions} de deux \textit{witnesses} sont placées côte à côte et reliées par de liens. Lorsque nous cliquons sur l'une d'entre elle, ils se colorent en rouge. Il est aussi possible d'afficher le score de similarité entre deux \textit{regions} en cochant la case \textit{Similarity score}. 

\subsubsection{Force-directed layout}

La visualisation \textit{Force-directed layout} représente les \textit{region pairs} sous forme de réseaux. 